\chapter{Scalar Field Sampling Along Trajectories}\label{ch:sampling}

Beyond advecting particles, the engine can record the value of model scalar fields
--- temperature and salinity, for example --- at each
particle position as the trajectory unfolds. This sampling is a diagnostic: it does
not affect the motion, but it attaches a time series of water properties to every
trajectory. The fields to sample are listed in the configuration key
\conf{mode.sample\_vars} (a single name is accepted as a one-element list; an empty list switches
sampling off and leaves the output byte-identical), and the horizontal method in
\conf{numerics.sample\_interp}. This chapter
describes what is sampled, where in the column the value is taken, how the horizontal
combination reuses the interpolation machinery of \cref{ch:hvel,ch:vertical}, and why
sampling is deliberately independent of the integration direction.

\section{What is sampled, and when}\label{sec:samp-what}

The sampled fields are node-based scalars, read from the same forcing dataset as the
velocity. Each requested field is probe-read once at startup, before the output file is created; a
field that is absent, or that is stored at element centres rather than nodes, raises a configuration
error immediately rather than failing midway through a run. Each reader prepares the fields in its own
way: the FESOM2 reader recognises mid-layer and interface fields by their vertical dimension and picks
the tracked level in a two-dimensional run; the FESOM-C reader reduces sigma-layered fields to two
dimensions with a node-aware reduction in a two-dimensional run; the HEALPix reader compacts pixels to nodes and fills each
column downward from its deepest valid value, so that the bottom clamp reads bottom water, never a
missing value (pixel geometry does not support sampling); the NEMO reader samples surface T-point
fields only. Values are recorded at the
trajectory output cadence (every \conf{output\_freq} steps, \cref{ch:output}), not at
every integration step. A record is labelled with the time at the end of its step (\cref{eq:out-time}), and
the scalar is taken at that instant as far as the forcing allows. When the velocity is reconstructed in time
(\cref{ch:temporal}), the scalar is re-interpolated to the exact record time, so the
sampled value is consistent with the recorded position; this re-interpolation is linear between
the bracketing snapshots even when the velocity uses the monotone-cubic reconstruction. In the
discrete-time mode the scalar is read from the snapshot at the end of the step --- the next snapshot, which the
step already holds as the target of its Runge--Kutta stages --- so its instant is the record time as well.
Only when no next snapshot exists, on the last step of a run that stops at the end of the forcing, is the
snapshot at the start of the step used, one snapshot behind the record time; the variable's comment
attribute says so. (Before output format~2 the discrete mode always sampled the start-of-step snapshot.)

Each sampled field becomes a variable of the trajectory output with the dimensions of the positions,
filled with missing values before its first record and in any record that does not carry it, and carrying
the units and long name of the source variable, the positions as its coordinates, and a comment on the
sampling instant; names that collide with the output's own variables, the release variables included, are
refused. Checkpoints do not
store sampled values. A field read internally for beaching (\cref{ch:lifecycle}) is written only when it
is also listed in \conf{sample\_vars}. On a rotated mesh (\cref{ch:frame}) scalars are sampled in the
rotated frame and written as they are, which is correct for scalars; a vector component listed as a
sample stays in the rotated frame.

\section{Vertical placement by field shape}\label{sec:samp-shape}

The vertical treatment is selected by the shape of the field array, which encodes
where the field lives in the column. A two-dimensional field, one value per node, is
sampled horizontally only. A mid-layer field, with one value per layer per node, is
interpolated linearly between the two bracketing mid-layer depths --- the midpoints of
the layers, computed from the fixed interface depths on a z-level mesh or from the
column geometry of \cref{eq:sigma} on a sigma mesh, and clamped to the first and last midpoint. On a
z-level mesh each neighbour is evaluated no deeper than its own floor, as for the vertical velocity
(\cref{sec:vert-whoriz}). On a sigma mesh each neighbour column is evaluated at the particle's
\emph{depth}, capped at the column thickness, and not along the particle's sigma surface as the
vertical velocity is; on a sloping bottom the sample therefore mixes neighbouring sigma levels, a
limitation that remains to be removed. An interface field, with one value per interface per node, is
handled by delegating to the vertical-velocity machinery of \cref{sec:vert-wprofile,sec:vert-whoriz},
including its sigma-following evaluation. In that delegation the integration-direction
factor is fixed to $+1$, for the reason given in \cref{sec:samp-direction}.

\section{Horizontal combination and element type}\label{sec:samp-horiz}

The horizontal combination offers the same four methods as the vertical velocity ---
\conf{nearest}, \conf{leastsq}, \conf{barycentric}, and \conf{idw} --- and shares their
behaviour with respect to the mesh. As with the vertical velocity, only the
\conf{barycentric} method depends on the element type, using the Wachspress weights of
\cref{sec:hvel-quad} on a quadrilateral and the triangular weights on a triangle; the
other three combine surrounding node values irrespective of element shape. Each method
falls back to the nearest node's value when its stencil is unusable. A particle that is not alive
(including a beached or a not-yet-released one) and a particle without a nearest node are recorded as
missing.

Model fields often carry fill values at coastal or below-bottom nodes, and a linear combination would
propagate them. The fitting methods therefore skip non-finite neighbours: the least-squares fit uses
the finite ones and needs at least three; inverse-distance weighting drops them; a barycentric element
with a non-finite vertex falls back to the nearest node's value. Under \conf{nearest}, and under
\conf{idw} for a particle that coincides with a node, the sample is that node's raw value, a fill value
included, because the sample point then \emph{is} that node; the least-squares fit has no such case and
returns the plane through the finite neighbours even at a node. Interface fields, which go through the vertical-velocity
machinery, are not filtered.

\section{Independence from the integration direction}\label{sec:samp-direction}

Sampling is a measurement of the instantaneous field and must not change sign when the
run is integrated backward in time. The engine enforces this by fixing the
direction factor to $+1$ in the sampling path, in contrast to the advection path, where
the factor $\tdir$ carries the forward or backward direction into the interpolated
velocity (\cref{sec:arch-direction}). The distinction matters specifically for an
interface-located scalar, which is sampled through the vertical-velocity machinery:
without the override it would inherit the velocity's direction sign and report a
sign-flipped value on a backward run. A sampled temperature is a temperature regardless
of the direction of integration.

\codeanchor{Sampling is implemented in \fn{sample\_field} and the direction-free combine
\fn{\_combine\_horizontal} (\srcfile{numerics/sampling.py}); the per-record sampling is
driven by \fn{\_sample\_record\_fields} (\srcfile{tracking.py}); each reader's
\fn{read\_scalar\_fields} prepares the fields and serves the startup probe.}
